%==============================================================================
% Language-Guided Optimization for Open-Vocabulary Gaussian Splatting SLAM
% RAL submission draft (IEEEtran journal format).
%
% NOTE: All quantitative result cells are rendered with \tofill (red "??").
%       They are PLACEHOLDERS and MUST be replaced with numbers obtained from
%       the ablation runs described in Sec. IV. Do not submit with placeholders.
%
% Compile:  pdflatex main && bibtex main && pdflatex main && pdflatex main
% Requires: IEEEtran.cls (RAL/IEEE template), references.bib
%==============================================================================
\documentclass[10pt, journal]{IEEEtran}

\usepackage{amsmath,amssymb,amsfonts}
\usepackage{graphicx}
\usepackage{booktabs}
\usepackage{multirow}
\usepackage{xcolor}
\usepackage{algorithm}
\usepackage{algpseudocode}
\usepackage{subcaption}
\usepackage[colorlinks=true,allcolors=blue]{hyperref}
\usepackage{url}

% Placeholder marker for results that must be filled from experiments.
\newcommand{\tofill}{\textcolor{red}{\textbf{??}}}
\newcommand{\best}[1]{\textbf{#1}}
\newcommand{\R}{\mathbb{R}}

\begin{document}

\title{Language-Guided Optimization for Open-Vocabulary \\ Gaussian Splatting SLAM}

\author{Anonymous Author(s)% <-- replace before camera-ready
\thanks{Manuscript received: Month, Day, Year; Revised Month, Day, Year; Accepted Month, Day, Year.}
\thanks{This paper was recommended for publication by Editor X upon evaluation of the Associate Editor and Reviewers' comments.}
\thanks{The authors are with the Department of XXX, University of XXX, City, Country (e-mail: author@univ.edu).}
\thanks{Digital Object Identifier (DOI): see top of this page.}}

\markboth{IEEE Robotics and Automation Letters. Preprint Version. Accepted Month, Year}{Author \MakeLowercase{\textit{et al.}}: Language-Guided Optimization for Open-Vocabulary Gaussian Splatting SLAM}

\maketitle

%==============================================================================
\begin{abstract}
Language-embedded 3D Gaussian Splatting (3DGS) SLAM produces photorealistic maps
that also support open-vocabulary queries, a capability increasingly required for
robotic interaction. In existing systems, however, the embedded language features
are treated as \emph{passive} attributes: they are distilled from a 2D
vision--language teacher and queried at test time, but they never inform the
geometric and photometric optimization of the map. We argue this decoupling
limits both rendering and segmentation quality. First, the distillation objective
(an $L_1$ loss) is misaligned with the cosine-similarity metric used for
open-vocabulary querying, and noisy 2D teacher predictions are baked into the 3D
field. Second, the map is densified only by uniform depth back-projection at
keyframes, so under-reconstructed and high-frequency/boundary regions are never
refined, capping rendering fidelity. We propose to make language an \emph{active}
optimization signal through four light-weight, online components: (i) a
CLIP-aligned distillation loss that matches the query metric; (ii) a per-pixel
uncertainty weighting derived from the entropy of a language codebook
soft-assignment; (iii) a geometry- and appearance-consistent feature smoothness
regularizer; and (iv) a language-boundary and joint-error guided densification
that spawns Gaussians where photometric error, geometric holes, or semantic
boundaries occur. All components reuse structures already present in a
language-embedded Gaussian SLAM and add negligible overhead. On Replica, ScanNet,
and TUM RGB-D, our method improves both rendering quality (PSNR) and
open-vocabulary segmentation (mIoU) over the baseline while maintaining real-time
operation. Code will be released.
\end{abstract}

\begin{IEEEkeywords}
SLAM, Gaussian Splatting, Open-Vocabulary Mapping, Semantic Scene Understanding, RGB-D Perception.
\end{IEEEkeywords}

%==============================================================================
\section{Introduction}\label{sec:intro}
\IEEEPARstart{D}{ense} visual SLAM that yields both photorealistic reconstruction
and semantic scene understanding is a long-standing goal for autonomous robots.
The recent emergence of 3D Gaussian Splatting (3DGS)~\cite{kerbl3Dgaussians} has
enabled SLAM systems that build high-fidelity, real-time renderable maps from
RGB-D or monocular video~\cite{keetha2024splatam,matsuki2024monogs,yan2024gsslam,yugay2023gaussianslam}.
In parallel, vision--language models such as CLIP~\cite{radford2021clip} and the
pixel-aligned LSeg~\cite{li2022lseg} have made it possible to embed
\emph{open-vocabulary} language features into 3D representations, allowing a robot
to localize arbitrary text queries (e.g., ``the mug on the table'') without a
closed label set~\cite{kerr2023lerf,qin2024langsplat,zhou2024feature3dgs}.
Combining the two lines of work, \emph{language-embedded Gaussian SLAM} attaches a
(compressed) language feature to every Gaussian and renders it differentiably,
producing maps that are simultaneously photorealistic and queryable.

Despite this progress, current language-embedded Gaussian SLAM systems treat the
language field as a \emph{passive} by-product. A 2D teacher (e.g., LSeg) produces
per-pixel features that are distilled into the per-Gaussian features and then
queried at test time, but the language signal is never fed back into the
optimization of the appearance and geometry of the map. We identify two concrete
consequences of this decoupling.

\textit{(1) The semantic objective is misaligned and noise-prone.} Distillation is
typically supervised with an $L_1$ loss on the raw feature vectors, whereas
open-vocabulary querying ranks classes by \emph{cosine} similarity to CLIP text
embeddings. This train/test mismatch caps the attainable segmentation accuracy.
Moreover, the 2D teacher is unreliable at object boundaries and in ambiguous
regions; an unweighted distillation loss bakes that noise directly into the 3D
field, further degrading the intersection-over-union (IoU).

\textit{(2) Densification ignores both error and semantics.} In RGB-D Gaussian
SLAM, new Gaussians are commonly created by uniformly back-projecting keyframe
depth. Without an error- or content-driven densification, under-reconstructed
regions, high-frequency texture, and thin structures never receive additional
Gaussians, imposing a ceiling on rendering fidelity (PSNR). Rendering-guided
densification has been explored for geometry~\cite{wu2024highfidelity}, but the
rich semantic signal already present in language-embedded maps remains unused for
this purpose.

We address both issues with a unifying idea: \textbf{make the language features an
active signal that drives the online optimization}. Concretely, we introduce four
light-weight components (Fig.~\ref{fig:teaser}). For semantic quality, we
(\textbf{A1}) align the distillation objective with the cosine query metric,
(\textbf{A2}) down-weight ambiguous pixels using the entropy of a language
codebook soft-assignment, and (\textbf{A3}) regularize the 3D field with a
geometry- and appearance-aware smoothness term. For rendering quality, we
(\textbf{B1}) densify the map at pixels that are simultaneously
under-reconstructed, photometrically erroneous, or located on a \emph{language
boundary}, which is exactly where additional Gaussians benefit both PSNR and
boundary IoU. All four components reuse structures already maintained by a
language-embedded Gaussian SLAM (the teacher features, the loop-closure codebook,
and the existing densification/pruning machinery), so they add negligible
overhead and preserve real-time operation.

Our contributions are:
\begin{itemize}
  \item A \emph{CLIP-aligned} feature distillation loss that matches the
  open-vocabulary cosine query metric, replacing the misaligned $L_1$ objective.
  \item An \emph{uncertainty-weighted} distillation scheme that suppresses noisy
  2D teacher supervision using codebook-entropy confidence, with no use of test
  labels.
  \item A \emph{geometry/appearance-consistent} feature smoothness regularizer
  that denoises the per-Gaussian language field.
  \item A \emph{language-boundary and joint-error guided} online densification
  that improves rendering fidelity and sharpens semantic boundaries.
  \item A unified, real-time system in which the above components consistently
  improve PSNR and mIoU over a strong language-embedded Gaussian SLAM baseline on
  Replica, ScanNet, and TUM RGB-D.
\end{itemize}

\begin{figure}[t]
  \centering
  \framebox[0.95\linewidth][c]{\rule{0pt}{3.2cm} \textit{[Teaser figure placeholder]}}
  \caption{Overview. Language features become an \emph{active} optimization signal:
  they (A1) align the distillation objective with the cosine query metric, (A2)
  weight supervision by semantic confidence, (A3) regularize the 3D field, and
  (B1) guide where the map is densified. \textbf{Replace with a qualitative
  figure (rendered RGB, segmentation, and Gaussian distribution near boundaries).}}
  \label{fig:teaser}
\end{figure}

%==============================================================================
\section{Related Work}\label{sec:related}

\subsection{Gaussian Splatting SLAM}
Following NeRF~\cite{mildenhall2020nerf} and NeRF-based dense
SLAM~\cite{zhu2022niceslam,wang2023coslam,sandstrom2023pointslam},
3DGS~\cite{kerbl3Dgaussians} has been adopted as an explicit, real-time renderable
map representation. SplaTAM~\cite{keetha2024splatam} performs RGB-D tracking and
mapping with silhouette-aware optimization; MonoGS~\cite{matsuki2024monogs}
introduces analytic camera-pose gradients through the rasterizer;
GS-SLAM~\cite{yan2024gsslam} and Gaussian-SLAM~\cite{yugay2023gaussianslam}
propose adaptive expansion and sub-map strategies;
Photo-SLAM~\cite{huang2024photoslam} targets efficient online operation. These
systems focus on geometry and appearance and do not embed open-vocabulary
semantics. Our work builds on a language-embedded RGB-D Gaussian SLAM and is
orthogonal to the tracking back-end.

\subsection{Semantic and Language-Embedded 3D Gaussians}
LERF~\cite{kerr2023lerf} distills CLIP features into a NeRF;
LangSplat~\cite{qin2024langsplat} and Feature-3DGS~\cite{zhou2024feature3dgs}
distill language/foundation features into 3D Gaussians for offline scenes. In the
SLAM setting, SGS-SLAM~\cite{li2024sgsslam} and SemGauss-SLAM~\cite{zhu2024semgauss}
attach \emph{closed-set} semantic labels to Gaussians and use them for mapping and
bundle adjustment. In contrast, we operate on \emph{open-vocabulary} language
features and, crucially, use them as an \emph{active} signal that shapes
densification and regularizes the map, rather than as a passive attribute to be
distilled and queried. Because high-dimensional features are expensive to store
per Gaussian, language-embedded SLAM compresses them (e.g., $512\!\to\!16$) with a
learned autoencoder; our objective-alignment and uncertainty terms are designed to
remain effective under this lossy compression.

\subsection{Adaptive Densification}
The original 3DGS densifies by cloning/splitting Gaussians with large
view-space positional gradients~\cite{kerbl3Dgaussians}. For SLAM,
rendering-guided densification has been used to map unobserved or re-observed
regions~\cite{wu2024highfidelity}, and RTG-SLAM~\cite{peng2024rtgslam} controls
Gaussian growth for scalability. These criteria are purely photometric/geometric.
We additionally exploit the \emph{semantic} structure of the scene: a language
boundary indicator steers densification toward object edges, which benefits both
rendering sharpness and boundary segmentation. To our knowledge, this is the first
densification criterion driven by an open-vocabulary language field.

%==============================================================================
\section{Method}\label{sec:method}

\subsection{System Overview and Preliminaries}\label{sec:prelim}
We build on a language-embedded RGB-D Gaussian SLAM with a tracking process and a
mapping process. Tracking estimates the camera pose by Generalized-ICP
(G-ICP)~\cite{segal2009gicp} between the incoming frame and the map. Mapping
maintains a set of anisotropic 3D Gaussians and optimizes them by differentiable
rasterization; loop closure uses a language bag-of-words over a codebook followed
by pose-graph optimization. We summarize only the elements needed for our method.

\textbf{Map representation.} The map is a set of Gaussians
$\{\mathcal{G}_i\}$, each parameterized by a mean $\boldsymbol{\mu}_i\!\in\!\R^3$,
a covariance $\Sigma_i$ (from scale $\mathbf{s}_i$ and rotation $\mathbf{q}_i$), an
opacity $o_i$, a view-dependent color $\mathbf{c}_i$, and a compressed language
feature $\mathbf{f}_i\!\in\!\R^{d}$ with $d\!=\!16$.

\textbf{Differentiable rendering.} For a pixel $\mathbf{p}$, color, depth, alpha,
and feature are rendered by front-to-back $\alpha$-blending of the $N$ Gaussians
overlapping $\mathbf{p}$:
\begin{equation}
\mathbf{X}(\mathbf{p}) = \sum_{i=1}^{N} \mathbf{x}_i\, \alpha_i \!\!\prod_{j<i}(1-\alpha_j),
\quad \mathbf{x}_i \in \{\mathbf{c}_i, z_i, 1, \mathbf{f}_i\},
\label{eq:render}
\end{equation}
yielding $\hat{C}$, $\hat{D}$, the accumulated opacity $A$, and the low-dimensional
feature map $\hat{F}\!\in\!\R^{d\times H\times W}$. A decoder
$\mathcal{D}_\psi\!:\!\R^{d}\!\to\!\R^{D}$ ($D\!=\!512$) lifts $\hat{F}$ to the
teacher's embedding space, and an encoder $\mathcal{E}_\phi\!:\!\R^{D}\!\to\!\R^{d}$
compresses teacher features when initializing Gaussians.

\textbf{Open-vocabulary query.} Given a set of text prompts encoded by CLIP into
$\{\mathbf{t}_k\}_{k=1}^{K}$, the predicted label at pixel $\mathbf{p}$ is
\begin{equation}
\hat{y}(\mathbf{p}) = \arg\max_{k}\; \cos\!\big(\mathcal{D}_\psi(\hat{F}(\mathbf{p})),\, \mathbf{t}_k\big).
\label{eq:query}
\end{equation}
Equation~\eqref{eq:query} is the basis of the mIoU/accuracy evaluation and
motivates the cosine alignment in Sec.~\ref{sec:a1}.

\textbf{Baseline objective.} The baseline optimizes
\begin{equation}
\mathcal{L}_{\text{base}} = (1\!-\!\lambda)\mathcal{L}_1^{\text{rgb}} + \lambda(1\!-\!\text{SSIM}) + \lambda_d \mathcal{L}_1^{\text{depth}} + \mathcal{L}_{\text{feat}},
\label{eq:base}
\end{equation}
with $\mathcal{L}_{\text{feat}}\!=\!\|\mathcal{D}_\psi(\hat{F})\!-\!\Phi\|_1$ and
$\Phi$ the 2D teacher (LSeg) feature. New Gaussians are added only by uniformly
back-projecting keyframe depth. Sections~\ref{sec:a1}--\ref{sec:b1} modify
$\mathcal{L}_{\text{feat}}$ and the densification.

\subsection{A1: CLIP-Aligned Feature Distillation}\label{sec:a1}
The query metric~\eqref{eq:query} depends only on cosine similarity, yet the
baseline supervises $\hat{F}$ with $L_1$. We add a channel-wise cosine term so the
training objective matches the metric. With valid mask
$\Omega=\{\mathbf{p}: \|\Phi(\mathbf{p})\|_1\!>\!0\}$ and per-pixel weight
$w(\mathbf{p})$ (Sec.~\ref{sec:a2}), and writing $\hat{\Phi}=\mathcal{D}_\psi(\hat{F})$,
\begin{equation}
\mathcal{L}_{\text{feat}} = \frac{\sum_{\mathbf{p}\in\Omega} w(\mathbf{p})\,\Big[ \ell_1(\mathbf{p}) + \lambda_{\cos}\big(1 - s_{\cos}(\mathbf{p})\big)\Big]}{\sum_{\mathbf{p}\in\Omega} w(\mathbf{p})},
\label{eq:a1}
\end{equation}
where $\ell_1(\mathbf{p})=\tfrac{1}{D}\|\hat{\Phi}(\mathbf{p})-\Phi(\mathbf{p})\|_1$
and $s_{\cos}(\mathbf{p})=\frac{\langle\hat{\Phi}(\mathbf{p}),\Phi(\mathbf{p})\rangle}{\|\hat{\Phi}(\mathbf{p})\|\,\|\Phi(\mathbf{p})\|}$
is computed over the channel dimension. Setting $\lambda_{\cos}\!=\!0$ and
$w\!\equiv\!1$ recovers the baseline, enabling a clean ablation.

\subsection{A2: Codebook-Entropy Uncertainty Weighting}\label{sec:a2}
The 2D teacher is least reliable at boundaries and in ambiguous regions. We derive
a per-pixel confidence from how decisively a teacher feature is assigned to the
scene language codebook $\mathcal{V}=\{\mathbf{v}_k\}_{k=1}^{K_c}$ ($K_c\!=\!64$),
which the SLAM system already builds for loop closure. With temperature $\tau$,
\begin{align}
\pi_k(\mathbf{p}) &= \frac{\exp(\cos(\hat{\Phi}_{\text{gt}}(\mathbf{p}),\mathbf{v}_k)/\tau)}{\sum_{k'}\exp(\cos(\hat{\Phi}_{\text{gt}}(\mathbf{p}),\mathbf{v}_{k'})/\tau)},\\
\tilde{H}(\mathbf{p}) &= -\tfrac{1}{\log K_c}\textstyle\sum_k \pi_k(\mathbf{p})\log \pi_k(\mathbf{p}),\\
w(\mathbf{p}) &= (1-\tilde{H}(\mathbf{p}))(1-w_{\min}) + w_{\min}.
\label{eq:a2}
\end{align}
A confident pixel (low entropy, dominated by one concept) gets $w\!\approx\!1$; an
ambiguous pixel is down-weighted to a floor $w_{\min}$ (we never discard a pixel
entirely). No text labels are used, so there is no evaluation leakage. The weight
enters~\eqref{eq:a1}.

\subsection{A3: Geometry/Appearance-Consistent Smoothness}\label{sec:a3}
Per-Gaussian features are supervised independently across views and remain noisy.
We add a bilateral $k$-NN consistency term that pulls together features of
Gaussians that are close in 3D space \emph{and} similar in color. For a random
subset $\mathcal{S}$ of Gaussians, with $\mathcal{N}_k(i)$ the $k$ nearest
neighbors of $i$ in $\mathcal{S}$,
\begin{align}
a_{ij} &= \exp\!\Big(\!-\tfrac{\|\boldsymbol{\mu}_i-\boldsymbol{\mu}_j\|^2}{\sigma_x^2}\Big)\exp\!\Big(\!-\tfrac{\|\mathbf{c}_i-\mathbf{c}_j\|^2}{\sigma_c^2}\Big),\\
\mathcal{L}_{\text{smooth}} &= \frac{\sum_{i\in\mathcal{S}}\sum_{j\in\mathcal{N}_k(i)} a_{ij}\big(1-\langle \hat{\mathbf{f}}_i,\hat{\mathbf{f}}_j\rangle\big)}{\sum_{i\in\mathcal{S}}\sum_{j\in\mathcal{N}_k(i)} a_{ij}},
\label{eq:a3}
\end{align}
where $\hat{\mathbf{f}}$ is the $\ell_2$-normalized feature and
$\boldsymbol{\mu},\mathbf{c}$ are detached so the term shapes only the features.
The affinity $a_{ij}$ deliberately weakens across color/geometry discontinuities,
so smoothing does not blur genuine semantic boundaries. Evaluated every
$T_{\text{sm}}$ iterations on $|\mathcal{S}|$ samples, the cost is bounded and
reuses the $k$-NN routine already used for language-based pruning.

\subsection{B1: Language-Boundary and Joint-Error Guided Densification}\label{sec:b1}
The baseline densifies only by uniform depth back-projection, so it cannot refine
regions that are visible but under-reconstructed, nor add resolution at semantic
edges. During mapping, when a keyframe $\mathcal{K}$ with pose $\mathbf{T}_{cw}$ is
rendered, we compute three per-pixel cues over valid-depth pixels:
\begin{align}
\text{photometric: } & e_{\text{rgb}}(\mathbf{p}) = \tfrac{1}{3}\|\hat{C}(\mathbf{p})-C_{\text{gt}}(\mathbf{p})\|_1,\\
\text{hole: } & h(\mathbf{p}) = \big[A(\mathbf{p})\!<\!\tau_\alpha\big] \vee \big[|\hat{D}(\mathbf{p})\!-\!D_{\text{gt}}(\mathbf{p})|\!>\!\tau_d\big],\\
\text{boundary: } & b(\mathbf{p}) = \big\|\nabla \hat{\Phi}_{\text{gt}}(\mathbf{p})\big\|\big/\textstyle\max_{\mathbf{p}'}\|\nabla \hat{\Phi}_{\text{gt}}(\mathbf{p}')\|,
\end{align}
where $b$ is the normalized magnitude of the spatial gradient of the normalized
teacher feature. The candidate set is
\begin{equation}
\mathcal{C} = \big\{\mathbf{p}: \text{valid}(\mathbf{p}) \wedge \big(h(\mathbf{p}) \vee e_{\text{rgb}}(\mathbf{p})\!>\!\tau_c \vee b(\mathbf{p})\!>\!\tau_b\big)\big\},
\end{equation}
ranked by a joint score $s(\mathbf{p})=e_{\text{rgb}}(\mathbf{p})+\gamma\, b(\mathbf{p})+\tau_c\,h(\mathbf{p})$.
We keep the top-$M$ candidates (and respect a global Gaussian cap), back-project
them with the camera intrinsics and $\mathbf{T}_{cw}$,
\begin{equation}
\boldsymbol{\mu}_{\text{new}} = \mathbf{R}_{cw}\Big[\tfrac{u-c_x}{f_x}z,\; \tfrac{v-c_y}{f_y}z,\; z\Big]^{\!\top} + \mathbf{t}_{cw},
\end{equation}
and spawn Gaussians with color from $C_{\text{gt}}$, feature
$\mathcal{E}_\phi(\Phi_{\text{gt}}(\mathbf{p}))$, identity rotation, opacity $0.1$,
and an isotropic scale equal to the metric pixel footprint $z/f_x$. New Gaussians
inherit the keyframe index so loop closure transforms them consistently. Because
densification targets exactly the high-error and boundary regions, it raises PSNR
while sharpening semantic edges; redundant Gaussians are removed by the existing
language-aware pruning.

\subsection{Total Objective and Online Schedule}
The mapping loss is
\begin{equation}
\mathcal{L} = (1\!-\!\lambda)\mathcal{L}_1^{\text{rgb}} + \lambda(1\!-\!\text{SSIM}) + \lambda_d \mathcal{L}_1^{\text{depth}} + \mathcal{L}_{\text{feat}} + \lambda_{\text{sm}}\mathcal{L}_{\text{smooth}},
\label{eq:total}
\end{equation}
with $\mathcal{L}_{\text{feat}}$ from~\eqref{eq:a1}. Algorithm~\ref{alg:main}
summarizes one mapping iteration; densification runs every $T_d$ iterations and
pruning every $T_p$, keeping per-frame cost approximately constant.

\begin{algorithm}[t]
\caption{One mapping iteration (additions in \textbf{bold}).}
\label{alg:main}
\begin{algorithmic}[1]
\State Select keyframe $\mathcal{K}$; render $\hat{C},\hat{D},A,\hat{F}$ via Eq.~\eqref{eq:render}
\State $\mathcal{L}_{\text{rgb,d}} \gets$ photometric + depth losses
\State \textbf{$w \gets$ codebook-entropy weights (Eq.~\eqref{eq:a2})} \Comment{A2}
\State \textbf{$\mathcal{L}_{\text{feat}} \gets$ CLIP-aligned distillation (Eq.~\eqref{eq:a1})} \Comment{A1}
\State \textbf{if} iter $\bmod\, T_{\text{sm}}=0$: \textbf{$\mathcal{L}\!\mathrel{+}=\!\lambda_{\text{sm}}\mathcal{L}_{\text{smooth}}$ (Eq.~\eqref{eq:a3})} \Comment{A3}
\State \textbf{if} iter $\bmod\, T_d=0$: \textbf{collect densify candidates $\mathcal{C}$} \Comment{B1}
\State $\mathcal{L}$.backward(); optimizer step on Gaussians and decoder
\State \textbf{if} candidates collected: \textbf{back-project top-$M$ and add Gaussians} \Comment{B1}
\State \textbf{if} iter $\bmod\, T_p=0$: language-aware pruning
\end{algorithmic}
\end{algorithm}

%==============================================================================
\section{Experiments}\label{sec:exp}

\subsection{Setup}
\textbf{Datasets.} We evaluate on Replica~\cite{straub2019replica} (8 sequences:
\textit{room0--2}, \textit{office0--4}), ScanNet~\cite{dai2017scannet} (6
sequences: \textit{0000, 0059, 0106, 0169, 0181, 0207}), and TUM
RGB-D~\cite{sturm2012tum} (\textit{fr1/desk}, \textit{fr2/xyz},
\textit{fr3/office}). Replica provides dense semantic ground truth for
open-vocabulary segmentation evaluation.

\textbf{Metrics.} Rendering: PSNR$\uparrow$, SSIM$\uparrow$,
LPIPS$\downarrow$. Open-vocabulary semantics: mIoU$\uparrow$ and pixel
accuracy$\uparrow$ via Eq.~\eqref{eq:query} using the dataset class names as text
prompts. Tracking: ATE-RMSE$\downarrow$ [cm]. Efficiency: tracking/mapping FPS,
final Gaussian count, and peak GPU memory.

\textbf{Baselines.} NICE-SLAM~\cite{zhu2022niceslam}, Co-SLAM~\cite{wang2023coslam},
Point-SLAM~\cite{sandstrom2023pointslam}, SplaTAM~\cite{keetha2024splatam},
MonoGS~\cite{matsuki2024monogs}, SGS-SLAM~\cite{li2024sgsslam}, and
SemGauss-SLAM~\cite{zhu2024semgauss}, plus the language-embedded Gaussian SLAM
\emph{baseline} we extend.

\textbf{Implementation.} We keep the baseline tracking/mapping unchanged and add
the four components. Unless noted, $\lambda_{\cos}\!=\!0.5$, $\tau\!=\!0.1$,
$w_{\min}\!=\!0.1$, $\lambda_{\text{sm}}\!=\!0.05$, $|\mathcal{S}|\!=\!2000$,
$k\!=\!8$, $T_{\text{sm}}\!=\!5$, $T_d\!=\!20$, $M\!=\!8000$, $\tau_\alpha\!=\!0.6$,
$\tau_d\!=\!5$\,cm, $\tau_c\!=\!0.08$, $\tau_b\!=\!0.5$, $\gamma\!=\!1.0$. All
experiments run on a single GPU. \emph{[Specify GPU model and report seeds.]}

\subsection{Rendering Quality}
Table~\ref{tab:render} reports rendering metrics on Replica. Our method improves
PSNR over the baseline, primarily from the language-boundary/error-guided
densification (B1) that refines under-reconstructed and high-frequency regions.

\begin{table}[t]
\centering
\caption{Rendering quality on Replica (averaged over 8 sequences). \textbf{Numbers are placeholders to be filled from experiments.}}
\label{tab:render}
\begin{tabular}{lccc}
\toprule
Method & PSNR$\uparrow$ & SSIM$\uparrow$ & LPIPS$\downarrow$ \\
\midrule
SplaTAM~\cite{keetha2024splatam} & \tofill & \tofill & \tofill \\
MonoGS~\cite{matsuki2024monogs}  & \tofill & \tofill & \tofill \\
SGS-SLAM~\cite{li2024sgsslam}    & \tofill & \tofill & \tofill \\
Baseline (language-embedded)     & \tofill & \tofill & \tofill \\
\textbf{Ours}                    & \best{\tofill} & \best{\tofill} & \best{\tofill} \\
\bottomrule
\end{tabular}
\end{table}

\subsection{Open-Vocabulary Semantic Segmentation}
Table~\ref{tab:sem} reports mIoU/accuracy on Replica. The CLIP-aligned objective
(A1) and uncertainty weighting (A2) contribute the largest semantic gains, with
A3 and B1 further sharpening boundaries.

\begin{table}[t]
\centering
\caption{Open-vocabulary 2D segmentation on Replica. \textbf{Placeholders.}}
\label{tab:sem}
\begin{tabular}{lcc}
\toprule
Method & mIoU$\uparrow$ & Acc.$\uparrow$ \\
\midrule
SGS-SLAM~\cite{li2024sgsslam} (closed-set)       & \tofill & \tofill \\
SemGauss-SLAM~\cite{zhu2024semgauss} (closed-set)& \tofill & \tofill \\
Baseline (open-vocab)                            & \tofill & \tofill \\
\textbf{Ours} (open-vocab)                       & \best{\tofill} & \best{\tofill} \\
\bottomrule
\end{tabular}
\end{table}

\subsection{Tracking Accuracy}
Our contributions target mapping; Table~\ref{tab:ate} verifies that tracking
accuracy is preserved (and may improve slightly as the refined map aids
frame-to-model alignment).

\begin{table}[t]
\centering
\caption{ATE-RMSE [cm] $\downarrow$. \textbf{Placeholders.}}
\label{tab:ate}
\begin{tabular}{lccc}
\toprule
Method & Replica & ScanNet & TUM \\
\midrule
Point-SLAM~\cite{sandstrom2023pointslam} & \tofill & \tofill & \tofill \\
SplaTAM~\cite{keetha2024splatam}         & \tofill & \tofill & \tofill \\
Baseline                                 & \tofill & \tofill & \tofill \\
\textbf{Ours}                            & \tofill & \tofill & \tofill \\
\bottomrule
\end{tabular}
\end{table}

\subsection{Ablation Study}
Table~\ref{tab:ablation} adds one component at a time on Replica, isolating each
contribution. We expect A1/A2/A3 to drive mIoU and B1 to drive PSNR, with the full
model best on both.

\begin{table}[t]
\centering
\caption{Component ablation on Replica. \textbf{Placeholders.}}
\label{tab:ablation}
\begin{tabular}{lcccc}
\toprule
Configuration & PSNR$\uparrow$ & mIoU$\uparrow$ & Acc.$\uparrow$ & FPS \\
\midrule
Baseline                       & \tofill & \tofill & \tofill & \tofill \\
\;+A1 (CLIP-aligned)           & \tofill & \tofill & \tofill & \tofill \\
\;+A1+A2 (uncertainty)         & \tofill & \tofill & \tofill & \tofill \\
\;+A1+A2+A3 (smoothness)       & \tofill & \tofill & \tofill & \tofill \\
\;+B1 (densification) = \textbf{Full} & \best{\tofill} & \best{\tofill} & \best{\tofill} & \tofill \\
\bottomrule
\end{tabular}
\end{table}

\subsection{Hyper-parameter Sensitivity}
We analyze sensitivity to $\lambda_{\cos}$, $\tau$, $\lambda_{\text{sm}}$, and the
densification budget $(T_d, M)$ in Table~\ref{tab:sens}, reporting the
PSNR/mIoU/FPS trade-off. \emph{[Fill from sweeps; recommend reporting at least 3
values per parameter.]}

\begin{table}[t]
\centering
\caption{Sensitivity (Replica, \textit{office0}). \textbf{Placeholders.}}
\label{tab:sens}
\begin{tabular}{lcccc}
\toprule
Parameter & Value & PSNR$\uparrow$ & mIoU$\uparrow$ & FPS \\
\midrule
\multirow{3}{*}{$\lambda_{\cos}$} & 0.25 & \tofill & \tofill & \tofill \\
                                  & 0.50 & \tofill & \tofill & \tofill \\
                                  & 1.00 & \tofill & \tofill & \tofill \\
\midrule
\multirow{3}{*}{$T_d$ (densify)}  & 10 & \tofill & \tofill & \tofill \\
                                  & 20 & \tofill & \tofill & \tofill \\
                                  & 50 & \tofill & \tofill & \tofill \\
\bottomrule
\end{tabular}
\end{table}

\subsection{Efficiency and Qualitative Results}
Table~\ref{tab:eff} reports runtime, Gaussian count, and memory; our components add
a small overhead while bounding map growth. Figure~\ref{fig:qual} shows that B1
concentrates Gaussians at object boundaries, yielding sharper renderings and
segmentation maps.

\begin{table}[t]
\centering
\caption{Efficiency on Replica. \textbf{Placeholders.}}
\label{tab:eff}
\begin{tabular}{lccc}
\toprule
Method & FPS & \#Gaussians (M) & Mem. (GB) \\
\midrule
Baseline      & \tofill & \tofill & \tofill \\
\textbf{Ours} & \tofill & \tofill & \tofill \\
\bottomrule
\end{tabular}
\end{table}

\begin{figure}[t]
  \centering
  \framebox[0.95\linewidth][c]{\rule{0pt}{3.0cm} \textit{[Qualitative comparison placeholder]}}
  \caption{Qualitative comparison (baseline vs.\ ours): rendered RGB, error map,
  open-vocabulary segmentation, and Gaussian density near boundaries.
  \textbf{Replace with real renders.}}
  \label{fig:qual}
\end{figure}

%==============================================================================
\section{Conclusion}\label{sec:conclusion}
We presented a language-guided optimization framework for open-vocabulary Gaussian
Splatting SLAM that turns the embedded language field from a passive attribute into
an active driver of the online map optimization. By aligning the distillation
objective with the cosine query metric, weighting supervision by semantic
confidence, regularizing the 3D field, and densifying at language boundaries and
high-error regions, our method improves both rendering fidelity and
open-vocabulary segmentation while preserving real-time operation, using only
structures already present in a language-embedded SLAM. Limitations include
reliance on the 2D teacher and the loop-closure codebook for the uncertainty term;
future work will explore self-supervised confidence and language-driven tracking.

%==============================================================================
\section*{Acknowledgments}
\emph{[Optional.]}

\bibliographystyle{IEEEtran}
\bibliography{references}

\end{document}
